\chapter{Verlet Laboratory and Energy Identity}
\label{app:integrator}

This appendix records the integrator and the energy identity used
in Chapter~15. It introduces no new physical hypothesis. The force
law is Assumption~\ref{ass:drag}. The bookkeeping reference is
Proposition~\ref{prop:allowed-drag}, which this laboratory does not
satisfy in the microscopic sense; the integrator tracks only
\(W_{\mathrm{diss}}\). The source file is
\texttt{code/namcg\_inspiral.py}.

\section{Reduced equation}
\label{sec:appC-ode}

In the reduced-mass frame the laboratory system is
\begin{equation}
  \label{eq:appC-ode}
  \ddot{\mathbf{r}}
  =\mathbf{a}_{\mathrm{N}}(\mathbf{r})
   +\mathbf{a}_{\mathrm{d}}(\mathbf{r},\mathbf{v}),
\end{equation}
with
\begin{equation}
  \label{eq:appC-acc}
  \mathbf{a}_{\mathrm{N}}
  =-\frac{G_{\mathrm{toy}}M}{d^{3}}\mathbf{r},
  \qquad
  \mathbf{a}_{\mathrm{d}}
  =-\frac{\eta}{d^{4}}\mathbf{v},
  \qquad
  d=|\mathbf{r}|.
\end{equation}
The conservative force is Newtonian. The drag is
\eqref{eq:Fdiss} with \(\mu\) cancelled. Toy parameters are
\eqref{eq:toy-params}.

\section{One step}
\label{sec:appC-step}

Let \(\Delta t>0\). Given \((\mathbf{r}_n,\mathbf{v}_n)\),
\begin{equation}
  \label{eq:appC-verlet}
  \begin{aligned}
    \mathbf{v}_{1/2}
    &=\mathbf{v}_n
      +\tfrac12\Delta t
       \bigl(
         \mathbf{a}_{\mathrm{N}}(\mathbf{r}_n)
         +\mathbf{a}_{\mathrm{d}}(\mathbf{r}_n,\mathbf{v}_n)
       \bigr),\\
    \mathbf{r}_{n+1}
    &=\mathbf{r}_n+\Delta t\,\mathbf{v}_{1/2},\\
    \mathbf{v}_{n+1}
    &=\mathbf{v}_{1/2}
      +\tfrac12\Delta t
       \bigl(
         \mathbf{a}_{\mathrm{N}}(\mathbf{r}_{n+1})
         +\mathbf{a}_{\mathrm{d}}(\mathbf{r}_{n+1},\mathbf{v}_{1/2})
       \bigr).
  \end{aligned}
\end{equation}
When \(\eta=0\) this is velocity Verlet for a position-dependent
force and is second-order accurate and symplectic on the
conservative phase space. When \(\eta\neq 0\) the last drag
evaluation uses \(\mathbf{v}_{1/2}\) rather than
\(\mathbf{v}_{n+1}\), so the dissipative piece is explicit and
the scheme is no longer symplectic. The local error in the drag
kick remains \(O(\Delta t^{2})\).

An implicit alternative replaces
\(\mathbf{a}_{\mathrm{d}}(\mathbf{r}_{n+1},\mathbf{v}_{1/2})\)
by \(\mathbf{a}_{\mathrm{d}}(\mathbf{r}_{n+1},\mathbf{v}_{n+1})\)
and solves a linear equation for \(\mathbf{v}_{n+1}\). The
published script does not do that. The difference is visible only
if one drives \(\Delta t\) up to the orbital period.

\section{Mechanical energy, work, and balance}
\label{sec:appC-energy}

Define
\begin{equation}
  \label{eq:appC-E}
  E_{\mathrm{mech}}
  =\frac12\mu v^{2}-\frac{G_{\mathrm{toy}}M\mu}{d}
\end{equation}
and the accumulated work
\begin{equation}
  \label{eq:appC-W}
  W_{\mathrm{diss}}(t_{n+1})
  =W_{\mathrm{diss}}(t_n)
   +\mathbf{F}_{\mathrm{diss}}(\mathbf{r}_n,\mathbf{v}_n)
    \cdot\mathbf{v}_n\,\Delta t,
\end{equation}
with \(\mathbf{F}_{\mathrm{diss}}=\mu\mathbf{a}_{\mathrm{d}}\) and
\(W_{\mathrm{diss}}(0)=0\). Along exact solutions of
\eqref{eq:appC-ode},
\begin{equation}
  \label{eq:appC-identity}
  \frac{\mathrm{d}E_{\mathrm{mech}}}{\mathrm{d}t}
  =\mathbf{F}_{\mathrm{diss}}\cdot\mathbf{v},
\end{equation}
because the Newtonian force is minus the gradient of
\(U_{\mathrm{N}}=-\Geff\)-less toy potential
\(-G_{\mathrm{toy}}M\mu/d\). Therefore
\begin{equation}
  \label{eq:appC-bal}
  \mathcal{E}(t)
  \equiv
  E_{\mathrm{mech}}(t)-W_{\mathrm{diss}}(t)
\end{equation}
is constant on exact solutions. On discrete solutions it is
constant up to the quadrature error in~\eqref{eq:appC-W} and the
Verlet error in~\eqref{eq:appC-verlet}. A systematic drift that
does not shrink with \(\Delta t\) means the implementation does
not obey~\eqref{eq:appC-identity} and is illegal in the sense of
Section~\ref{sec:ch14-legal}.

A trapezoidal quadrature,
\begin{equation}
  \label{eq:appC-trap}
  \Delta W
  =\frac12
   \Bigl(
     \mathbf{F}_{\mathrm{diss},n}\cdot\mathbf{v}_n
     +\mathbf{F}_{\mathrm{diss},n+1}\cdot\mathbf{v}_{n+1}
   \Bigr)\Delta t,
\end{equation}
reduces that drift by one order in \(\Delta t\). The published
script uses the left-endpoint rule~\eqref{eq:appC-W} because it
is enough to make \(\mathcal{E}\) look flat on the scale of
Figure~\ref{fig:lab}.

\section{Angular momentum}
\label{sec:appC-L}

The Newtonian force is central and exerts no torque. The drag is
antiparallel to \(\mathbf{v}\), so
\begin{equation}
  \label{eq:appC-torque}
  \frac{\mathrm{d}\mathbf{L}}{\mathrm{d}t}
  =\mathbf{r}\times\mathbf{F}_{\mathrm{diss}}
  =-\eta\mu\frac{\mathbf{r}\times\mathbf{v}}{d^{4}}.
\end{equation}
Angular momentum is therefore not conserved when \(\eta\neq 0\).
That is required: an inspiral that conserved \(|\mathbf{L}|\) at
fixed mass would be incompatible with a shrinking circular orbit.
The identity~\eqref{eq:appC-torque} is a useful code check
(Exercise~4 of Chapter~15).

\section{Initial data and units}
\label{sec:appC-ic}

The published run uses
\begin{equation}
  \label{eq:appC-ic}
  \mathbf{r}(0)=(1,0),\qquad
  \mathbf{v}(0)=(0,\sqrt{G_{\mathrm{toy}}M}),
\end{equation}
with \(G_{\mathrm{toy}}=36\), \(\eta=0.006\), \(\Delta t=10^{-3}\),
and \(1.2\times 10^{4}\) steps. These are code units. They are
chosen so that decay is visible, not so that \(\eta\) matches a
tidal-disruption estimate from Problem~I of Chapter~17.

Restoring dimensions, \(\eta\) in~\eqref{eq:appC-acc} has dimension
\(\mathrm{length}^{4}/\mathrm{time}\), since
\(\mathbf{a}=-\eta\mathbf{v}/d^{4}\). The static tail
\eqref{eq:I-tail} does not supply this coefficient; deriving it
requires the response calculation of Problem~\ref{sec:status-I},
which this appendix does not attempt.

\section{How to run the script}
\label{sec:appC-run}

From the project root, after \texttt{figs/} exists:
\begin{verbatim}
python code/namcg_inspiral.py
\end{verbatim}
The script writes
\texttt{figs/namcg\_gr\_simulation\_corrected.pdf}. Point
Figure~\ref{fig:lab} at that PDF for the print interior. A
legacy PNG from an earlier notebook may be kept for draft
viewing.

If \texttt{matplotlib} is unavailable, the integrator still runs
and can write \texttt{traj}, \texttt{E\_mech}, and
\texttt{W\_diss} to text; the figure is then optional.

\section{What this appendix will not do}
\label{sec:appC-not}

It will not introduce gravitational waves, horizons, or a
relativistic two-body Hamiltonian. It will not treat
\(G_{\mathrm{toy}}\) as \(G_{\mathrm{N}}\). It will not fit
\(\eta\) to data. Those refusals are the content of
Sections~\ref{sec:ch15-misread} and~\ref{sec:ch16-public}.

\section{Exercises}
\label{sec:appC-exercises}

\begin{enumerate}
  \item \textbf{Conservative order.}
    With \(\eta=0\), measure the secular drift of
    \(E_{\mathrm{mech}}\) over \(10^{4}\) steps at three values
    of \(\Delta t\). Confirm second-order scaling.
  \item \textbf{Quadrature.}
    Replace~\eqref{eq:appC-W} by~\eqref{eq:appC-trap} and compare
    the drift of \(\mathcal{E}\).
  \item \textbf{Implicit drag.}
    Solve the linear equation for \(\mathbf{v}_{n+1}\) when the
    last drag evaluation uses \(\mathbf{v}_{n+1}\). Does
    \(\mathcal{E}\) improve enough to justify the extra algebra?
  \item \textbf{Torque check.}
    Integrate~\eqref{eq:appC-torque} in lockstep with the
    trajectory and compare with \(\mathbf{L}_{n+1}-\mathbf{L}_n\).
  \item \textbf{Illegal control.}
    Set \(\mathbf{a}_{\mathrm{d}}=-\gamma\mathbf{v}\) with
    constant \(\gamma\) and watch \(W_{\mathrm{diss}}\) as
    \(d\) decreases. This is Exercise~3 of Chapter~15, recorded
    here so the appendix can be used alone as a lab sheet.
\end{enumerate}